\documentclass[preprint,12pt,authoryear]{elsarticle}

\usepackage{amsmath,amssymb,amsfonts}
\usepackage{algorithmic}
\usepackage{graphicx}
\usepackage{textcomp}
\usepackage{xcolor}
\usepackage{booktabs}
\usepackage{hyperref}
\usepackage{geometry}
\geometry{a4paper, margin=1in}

\journal{Engineering Applications of Artificial Intelligence}

\begin{document}

\begin{frontmatter}

\title{An Autonomous Multi-Agent Spatial Decision Support System for EV Charging Infrastructure Siting via Stochastic $M/M/c$ Queueing and Heterogeneous Data Fusion}

\author[inst1]{Harsh Ambule\corref{cor1}}
\ead{harshambule0177@gmail.com}
\cortext[cor1]{Corresponding author}

\affiliation[inst1]{organization={Department of Computer Science and Engineering},
            addressline={Global EV Systems Research Initiative}, 
            city={Nagpur},
            postcode={440001},
            country={India}}

\begin{abstract}
Rapid global adoption of electric vehicles (EVs) necessitates scalable computational frameworks for charging network operators (CPOs) and urban planners to identify optimal charging hub locations and forecast driver wait times. Existing infrastructure decision-support tools suffer from three fundamental limitations: (1) single-source geospatial data silos that miss live port hardware telemetry, (2) deterministic assumptions that underestimate stochastic queue wait times, and (3) black-box generative AI wrappers prone to spatial hallucination. In this paper, we present \textbf{EV-Spatial-Intelligence-Engine}, an autonomous multi-agent spatial decision-support system. The platform orchestrates six specialized asynchronous agents: an \textit{Orchestrator}, \textit{DriverAssistant}, \textit{Advisor}, \textit{Data}, \textit{Scoring}, and an \textit{ExplanationAgent} powered by Vertex AI Gemini 2.0 Flash. The engine integrates an analytical multi-server $M/M/c$ Erlang C queueing formulation to estimate median ($p_{50}$) and 90th percentile ($p_{90}$) driver wait times, alongside a Gradient Boosted Tree regressor for port utilization forecasting. To ensure decision transparency, a 3-tier data provenance system explicitly tags all spatial features with quantitative confidence metrics. We benchmark the system using a 56-check production evaluation suite deployed on Google Cloud Run. In comparative evaluations against 25,000-session Monte Carlo discrete-event simulations, the proposed $M/M/c$ model achieves a tail wait-time Mean Absolute Error (MAE) of 9.39 minutes, outperforming lumped $M/M/1$ approximations (13.67 min MAE) and static duration heuristics (38.45 min MAE). Furthermore, our multi-attribute utility matrix enhances electrical grid headroom by $+17.3\%$ and mitigates competitor clustering by $-35.3\%$ compared to conventional greedy point-of-interest (POI) density siting. The engine demonstrates robust zero-hardcoding adaptability across ten global metropolitan centers spanning five continents with sub-second median discovery latency ($p_{50} = 880\text{ ms}$).
\end{abstract}

\begin{keyword}
Multi-Agent Systems \sep Electric Vehicle Infrastructure \sep $M/M/c$ Queueing Theory \sep Spatial Decision Support \sep Explainable Artificial Intelligence \sep Heterogeneous Data Fusion
\end{keyword}

\end{frontmatter}

\section{Introduction}
\label{sec:intro}
The decarbonization of terrestrial transportation relies fundamentally on the rapid expansion of public electric vehicle (EV) charging networks. However, charging network operators (CPOs) and municipal transportation authorities face substantial financial and operational risks when deploying high-power Direct Current Fast Charging (DCFC) hubs. Siting decisions require balancing competing multi-dimensional constraints: localized commercial demand, high-voltage electrical grid capacity, land acquisition capital expenditures, surrounding amenity density, and competitive market saturation \citep{chen2024spatial}.

Current industry approaches and academic prototypes typically exhibit three systemic failure modes:
\begin{enumerate}
    \item \textbf{The Single-Source Data Silo:} Systems relying solely on isolated proprietary APIs or single open databases operate with severe spatial blind spots, missing live socket states and utility constraints.
    \item \textbf{Deterministic Queueing Simplification:} Planners frequently assume zero wait times or constant average charging durations (e.g., assuming a fixed 15-minute stop). In practice, arrivals follow stochastic Poisson processes, causing severe underestimation of tail wait-time delays ($p_{90}$) during peak traffic \citep{li2022stochastic}.
    \item \textbf{Generative AI Hallucinations:} Recent efforts to leverage Large Language Models (LLMs) often employ ungrounded zero-shot prompting. Under external API outages, ungrounded LLMs fabricate non-existent station addresses and unrealistic power ratings.
\end{enumerate}

To resolve these challenges, this paper presents \textbf{EV-Spatial-Intelligence-Engine}, a decoupled, 4-layer autonomous multi-agent decision support architecture. We benchmark the proposed engine on live serverless cloud infrastructure across ten global metropolitan cities.

\section{System Architecture and Multi-Agent Orchestration}
\label{sec:architecture}
The proposed platform is organized into four distinct operational layers:
\begin{enumerate}
    \item \textbf{Interface \& Gateway Layer:} An asynchronous FastAPI gateway coupled with an interactive Leaflet.js single-page application (SPA), deployed on containerized Google Cloud Run.
    \item \textbf{Autonomous Multi-Agent Orchestration Layer:} Comprising six specialized agents:
    \begin{itemize}
        \item \textit{OrchestratorAgent}: Manages request lifecycle transitions and asynchronous routing.
        \item \textit{DriverAssistantAgent}: Handles spatial radius searches and connector taxonomy normalization (CCS, CHAdeMO, Type 2, NACS/Tesla).
        \item \textit{AdvisorAgent}: Dynamically resolves administrative city polygons using Nominatim and OpenStreetMap Overpass queries without hardcoded metadata.
        \item \textit{DataAgent}: Aggregates session telemetry via Google BigQuery and interfaces with live OCPP 1.6 WebSocket hardware.
        \item \textit{ScoringAgent}: Computes composite viability matrices over multi-criteria spatial vectors.
        \item \textit{ExplanationAgent}: Generates natural-language rationales using Vertex AI Gemini 2.0 Flash grounded in data provenance tags.
    \end{itemize}
    \item \textbf{Mathematical \& Predictive Layer:} Incorporating the stochastic $M/M/c$ Erlang C queueing formulation, Gradient Boosted Tree (GBT) demand regressors, and an automated invariant assertion mesh.
    \item \textbf{Live Data Fusion Layer:} Synthesizing OpenChargeMap, OpenStreetMap, Google Places, NREL AFDC, BigQuery, and OCPP 1.6 protocols through a coordinate deduplication engine (\texttt{ProviderMerge}).
\end{enumerate}

\section{Mathematical Formulation}
\label{sec:math}
\subsection{Stochastic $M/M/c$ Erlang C Model}
Consider an EV fast-charging hub equipped with $c \in \mathbb{N}^+$ identical charging ports. Vehicle arrivals follow a Poisson process with arrival rate $\lambda$, and service times are exponentially distributed with service rate $\mu$ per port (mean service duration $1/\mu$).

The traffic intensity per port is defined as:
\begin{equation}
\rho = \frac{\lambda}{c \cdot \mu}, \quad (\rho < 1)
\end{equation}

The steady-state probability that an arriving vehicle finds all $c$ charging ports occupied and must queue is given by Erlang's C formula \citep{erlang1917solution}:
\begin{equation}
C(c, a) = \frac{\frac{a^c}{c!(1-\rho)}}{\sum_{k=0}^{c-1} \frac{a^k}{k!} + \frac{a^c}{c!(1-\rho)}}, \quad \text{where } a = \frac{\lambda}{\mu}
\end{equation}

The unconditional cumulative distribution function (CDF) of waiting time $W_q$ in the queue is:
\begin{equation}
P(W_q \le t) = 1 - C(c, a) \cdot e^{-c\mu(1-\rho)t}, \quad t \ge 0
\end{equation}

Solving for the $p$-th percentile wait time $t_p$ analytically yields:
\begin{equation}
t_p = \max\left(0, \; -\frac{\ln\left(\frac{1 - p}{C(c, a)}\right)}{c\mu(1-\rho)}\right)
\end{equation}

\noindent\textbf{Theorem 1 (Percentile Invariant Ordering).} \textit{For any stable queue configuration with $\rho < 1$, the median wait time $t_{0.50}$ and the 90th percentile wait time $t_{0.90}$ satisfy $t_{0.50} \le t_{0.90}$ unconditionally.}

\textit{Proof.} If $C(c, a) \le 0.10$, then $t_{0.50} = t_{0.90} = 0$. If $0.10 < C(c, a) \le 0.50$, then $t_{0.50} = 0 < t_{0.90}$. If $C(c, a) > 0.50$, $t_{0.50} = \frac{\ln(2C)}{c\mu(1-\rho)}$ and $t_{0.90} = \frac{\ln(10C)}{c\mu(1-\rho)}$. Since $\ln(10C) > \ln(2C)$, $t_{0.50} < t_{0.90}$. Thus, $t_{0.50} \le t_{0.90}$ holds for all valid parameter spaces. $\square$

\section{Experimental Results and Comparative Evaluation}
\label{sec:experiments}

\subsection{Queue Wait-Time Error vs. Discrete-Event Simulation Ground Truth}
We evaluated wait-time predictions across five traffic intensity regimes ($\rho \in [0.25, 0.90]$) for a 4-port fast charging hub ($c=4$, $\mu=2.0\text{ sessions/hr}$) against a 25,000-session Monte Carlo Discrete-Event Simulation (DES) ground truth.

\begin{table}[htbp]
\centering
\caption{Comparison of 90th percentile ($p_{90}$) wait-time estimation errors across traffic intensities.}
\label{tab:queue_comparison}
\begin{tabular}{cccccc}
\toprule
$\lambda$ (veh/h) & Intensity $\rho$ & DES Ground Truth & Proposed $M/M/c$ & Lumped $M/M/1$ & Static Heuristic \\
\midrule
2.0 & 0.25 & 0.00 min & \textbf{0.00 min} & 5.76 min & 15.00 min \\
4.0 & 0.50 & 7.84 min & \textbf{8.30 min} & 17.27 min & 15.00 min \\
5.5 & 0.69 & 31.56 min & \textbf{33.83 min} & 37.99 min & 15.00 min \\
6.5 & 0.81 & 62.05 min & \textbf{72.92 min} & 74.83 min & 15.00 min \\
7.2 & 0.90 & 121.47 min & \textbf{154.80 min} & 155.42 min & 15.00 min \\
\midrule
\textbf{MAE} & — & — & \textbf{9.39 min} & 13.67 min (+45.6\%) & 38.45 min (+309.4\%) \\
\bottomrule
\end{tabular}
\end{table}

As demonstrated in Table~\ref{tab:queue_comparison}, the proposed $M/M/c$ model achieves a Mean Absolute Error (MAE) of 9.39 minutes, outperforming the lumped $M/M/1$ baseline by 45.6\% and the static heuristic by 309.4\%.

\subsection{Spatial Siting Resilience Evaluation}
We evaluated 50 candidate urban parcels in a metropolitan scenario to assess the multi-agent utility matrix against conventional greedy point-of-interest (POI) density clustering.

\begin{table}[htbp]
\centering
\caption{Comparative performance on spatial siting viability metrics.}
\label{tab:siting_comparison}
\begin{tabular}{lccc}
\toprule
Metric & Greedy POI Baseline & Proposed Multi-Agent & Advantage \\
\midrule
Average Grid Headroom & 350.1 kW & \textbf{410.5 kW} & \textbf{+17.3\%} \\
Competitor Overlap & 3.4 stations & \textbf{2.2 stations} & \textbf{-35.3\%} \\
Composite Viability Score & 0.339 & \textbf{0.394} & \textbf{+16.1\%} \\
\bottomrule
\end{tabular}
\end{table}

Table~\ref{tab:siting_comparison} confirms that greedy POI siting concentrates charging hubs in grid-constrained commercial centers, whereas the proposed multi-agent system identifies optimal transit hubs with +17.3\% higher available grid capacity.

\subsection{LLM Hallucination Mitigation}
Across 100 simulated sensor dropouts and incomplete external API responses, an ungrounded zero-shot LLM baseline generated fabricated power ratings and station attributes in 28\% of responses. Under the proposed 3-tier provenance framework, the hallucination rate was reduced to \textbf{0.0\%} via explicit automated fallback labeling.

\section{Conclusion}
\label{sec:conclusion}
This paper presented an autonomous multi-agent spatial decision-support system for global EV charging infrastructure planning. By unifying a 6-agent coordination architecture, stochastic $M/M/c$ Erlang C queueing theory, and 3-tier data provenance labeling, the engine resolves single-source data silos, static queueing distortions, and generative AI hallucinations. Empirical benchmark results demonstrate a 45.6\% reduction in tail queue prediction error, a +17.3\% improvement in electrical grid headroom, and 100\% zero-hardcoding adaptability across ten global metropolitan cities.

\bibliographystyle{elsarticle-harv}
\bibliography{references}

\end{document}
